\textbf{واحد:} 2 واحد (1 نظری + 1 عملی)، در مجموع 48 ساعت · \textbf{شکل
برگزاری:} هر هفته 1 ساعت درس + 2 ساعت \lr{lab} \textbf{مخاطبان:}
دانشجویان کارشناسی ارشد و سال آخر کارشناسی · \textbf{پیش‌نیازها:} روش‌های
آماری در علوم شناختی؛

\subsection{اصول طراحی}

\begin{itemize}
\tightlist
\item
  \textbf{یک مسیر مدل‌سازی ثابت برای کل ترم.} همه‌ی \lr{lab}ها همین مراحل
  را دنبال می‌کنند: \emph{شبیه‌سازی \lr{(simulate)} \ensuremath{\leftarrow} برازش \lr{(fit)} \ensuremath{\leftarrow}
  ارزیابی \lr{(evaluate)} \ensuremath{\leftarrow} بازیابی \lr{(recover)} \ensuremath{\leftarrow} تفسیر
  \lr{(interpret)}}.
\item
  \textbf{چند مجموعه‌داده‌ی ثابت که مدام تکرار می‌شوند.} منحنی‌های فراموشی و
  نظریه‌ی تشخیص علامت \lr{(SDT)} از هفته‌ی 2 تا 13 ادامه دارند، تا روش‌های
  جدید روی داده‌های آشنا مقایسه شوند. فصل 6 دسته‌بندی \lr{(GCM)}، زمان
  واکنش \lr{(drift-diffusion)} و یادگیری تقویتی \lr{(reinforcement
  learning)} را اضافه می‌کند.
\item
  \textbf{وزن فصل‌ها.} فصل 1: 2 هفته · فصل 2: 3 · فصل 3: 2 · فصل 4: 2 ·
  فصل 5: 3 · فصل 6: 2. سه هفته‌ی دیگر هم به معارفه، میان‌ترم و ارائه‌ی
  پروژه‌ها اختصاص دارد.
\item
  \textbf{\lr{«Metropolis\lr{--}Gibbs»} یعنی دو الگوریتم جدا:}
  \lr{Metropolis\lr{--}Hastings} (هفته‌ی 11) و نمونه‌گیری \lr{Gibbs} (هفته‌ی
  12). روش \lr{MH-within-Gibbs} در یک کادر «پیشرفته» آمده است.
\end{itemize}

\subsection{برنامه‌ی هفتگی}

\begin{longtable}[]{@{}
  >{\raggedright\arraybackslash}p{(\linewidth - 12\tabcolsep) * \real{0.1429}}
  >{\raggedright\arraybackslash}p{(\linewidth - 12\tabcolsep) * \real{0.1429}}
  >{\raggedright\arraybackslash}p{(\linewidth - 12\tabcolsep) * \real{0.1429}}
  >{\raggedright\arraybackslash}p{(\linewidth - 12\tabcolsep) * \real{0.1429}}
  >{\raggedright\arraybackslash}p{(\linewidth - 12\tabcolsep) * \real{0.1429}}
  >{\raggedright\arraybackslash}p{(\linewidth - 12\tabcolsep) * \real{0.1429}}
  >{\raggedright\arraybackslash}p{(\linewidth - 12\tabcolsep) * \real{0.1429}}@{}}
\caption{\centering برنامه‌ی هفتگی}\tabularnewline
\toprule\noalign{}
\begin{minipage}[b]{\linewidth}\raggedright
هفته
\end{minipage} & \begin{minipage}[b]{\linewidth}\raggedright
فصل
\end{minipage} & \begin{minipage}[b]{\linewidth}\raggedright
درس (1 ساعت)
\end{minipage} & \begin{minipage}[b]{\linewidth}\raggedright
\lr{lab} (2 ساعت)
\end{minipage} & \begin{minipage}[b]{\linewidth}\raggedright
اهداف یادگیری
\end{minipage} & \begin{minipage}[b]{\linewidth}\raggedright
منابع
\end{minipage} & \begin{minipage}[b]{\linewidth}\raggedright
تکلیف / نقطه‌ی عطف
\end{minipage} \\
\midrule\noalign{}
\endfirsthead
\toprule\noalign{}
\begin{minipage}[b]{\linewidth}\raggedright
هفته
\end{minipage} & \begin{minipage}[b]{\linewidth}\raggedright
فصل
\end{minipage} & \begin{minipage}[b]{\linewidth}\raggedright
درس (1 ساعت)
\end{minipage} & \begin{minipage}[b]{\linewidth}\raggedright
\lr{lab} (2 ساعت)
\end{minipage} & \begin{minipage}[b]{\linewidth}\raggedright
اهداف یادگیری
\end{minipage} & \begin{minipage}[b]{\linewidth}\raggedright
منابع
\end{minipage} & \begin{minipage}[b]{\linewidth}\raggedright
تکلیف / نقطه‌ی عطف
\end{minipage} \\
\midrule\noalign{}
\endhead
\bottomrule\noalign{}
\endlastfoot
1 & 1 & \textbf{معارفه: مدل شناختی چیست؟} چرا مدل می‌سازیم؛ نظریه‌ی کلامی
در برابر فرمال؛ سه سطح \lr{Marr}؛ خانواده‌های مدل (نمادین، پیوندگرا یا
\lr{connectionist}، تصادفی/پویا، بیزی)؛ نقشه‌ی درس؛ \textbf{جدول
نمادگذاری} & شبیه‌سازی گام تصادفی و مقایسه‌ی پیش‌بینی‌های مدل با مشاهده‌ها &
تمایز مدل‌های توصیفی، فرایندی و هنجاری؛ توضیح اینکه چرا نظریه‌های کلامی
ناقص‌اند؛ اجرا و رسم شبیه‌سازی یک مدل & \lr{F\lr{\&}L} فصل 1\lr{--}2؛ \lr{Verguts}،
فصل مقدمه {[}بررسی شود{]}؛ \lr{Anderson}، فصل مقدمه {[}بررسی شود{]} &
\lr{A0 (}بدون نمره): یادداشتی یک‌صفحه‌ای: «یک نظریه‌ی کلامی انتخاب کنید و بنویسید چه چیزهایی را مشخص نکرده است» \\
2 & 1 & \textbf{کالبدشکافی یک مدل:} داده، پارامتر آزاد در برابر ثابت،
معادلات مدل، مدل نویز \lr{(noise/data model)}؛ شبیه‌سازی به‌عنوان ابزار
فکر کردن؛ اولین مدل‌ها: \lr{SDT} (\(d'\)، \(c\)) و فراموشی نمایی در برابر
توانی & صورت‌بندی \lr{SDT} (\(d'\)، \(c\)، منحنی \lr{ROC)}؛ شبیه‌سازی
فراموشی نمایی و توانی؛ بررسی اثر هر پارامتر روی پیش‌بینی‌ها & شناسایی
اجزای یک مدل فرمال؛ محاسبه‌ی \(d'\) و \(c\)؛ پیش‌بینی اثر پارامترها روی
خروجی مدل & \lr{F\lr{\&}L} فصل 2؛ \lr{L\lr{\&}W} فصل 11 (پیش‌نمایش \lr{SDT)} &
\textbf{انتشار \lr{A1:}} «از کلامی به فرمال» (موعد: هفته‌ی 4) \\
3 & 2 & \textbf{کمترین مربعات و بهینه‌سازی:} \lr{SSE}، \lr{MSE}،
\lr{RMSD}؛ سطح خطا؛ جستجوی شبکه‌ای، \lr{Nelder\lr{--}Mead}، روش‌های گرادیانی؛
کمینه‌ی محلی، مقدار شروع، کران‌ها و \lr{reparameterization} &
برازش منحنی فراموشی با جستجوی شبکه‌ای ؛ رسم سطح خطا؛ برازش چندشروعی (\lr{multi-start}
&
استخراج و ارزیابی برازش \lr{LS}؛ مقایسه‌ی بهینه‌سازها؛ تشخیص برازش‌های ناموفق
& \lr{F\lr{\&}L} فصل 3؛ قدوسیان (روش‌های جستجوی مستقیم) & \lr{---} \\
4 & 2 & \textbf{بیشینه درست‌نمایی:} احتمال در برابر درست‌نمایی؛
\lr{log-likelihood}؛ درست‌نمایی برنولی، دوجمله‌ای و گاوسی؛ \lr{LS} به‌عنوان
\lr{MLE} تحت نویز گاوسی & \lr{MLE} برای داده‌ی فراموشی (تعداد یادآوری
دوجمله‌ای) و برای \lr{SDT}؛ مقایسه‌ی برآوردهای \lr{MLE} و \lr{LS} &
استخراج درست‌نمایی؛ تحلیل کمینه‌سازی \lr{NLL}؛ توضیح اینکه \lr{LS }و \lr{MLE }کی یکسان‌اند
& \lr{F\lr{\&}L} فصل 4؛ عظیمی‌فر و احمدی (برآورد \lr{ML)} & \textbf{موعد
\lr{A1} · انتشار \lr{A2:}} \lr{LS} در برابر \lr{MLE} (موعد: هفته‌ی 6) \\
5 & 2 & \textbf{عدم قطعیت و چند شرکت‌کننده:} خطای استاندارد از
\lr{Hessian}، \lr{profile likelihood}، \lr{bootstrap} پارامتری؛ برازش
فردی در برابر تجمیعی؛ خطاهای ناشی از میانگین‌گیری & بازه‌ی اطمینان
\lr{bootstrap}؛ میانگین‌گیری از یادگیرنده‌های نمایی یک «قانون توانی تمرین»
کاذب می‌سازد؛ برازش‌های فردی & سنجش عدم قطعیت برآورد؛ توضیح خطای
میانگین‌گیری؛ انتخاب سطح برازش & \lr{F\lr{\&}L} فصل 5 & \lr{---} \\
6 & 3 & \textbf{ارزیابی 1، برازش در برابر تعمیم‌پذیری:} \(R^2\)،
\lr{RMSD} و \lr{deviance} و محدودیت‌هایشان؛ انعطاف‌پذیری و
\lr{overfitting}؛ اعتبارسنجی \lr{hold-out}، \lr{k-fold} و \lr{LOO} &
نمایش \lr{overfitting} با چندجمله‌ای؛ مقایسه‌ی سه مدل فراموشی با
\lr{cross-validation} &
توضیح اینکه چرا برازش خوب \ensuremath{\neq} مدل خوب؛ تحلیل \lr{k-fold CV}؛ تفسیر نتایج \lr{CV}
& \lr{F\lr{\&}L} فصل 10؛ عظیمی‌فر و احمدی (تعمیم‌پذیری) & \textbf{موعد \lr{A2}
· انتشار \lr{A3:}} مقایسه + بازیابی (موعد: هفته‌ی 9) \\
7 & 3 & \textbf{ارزیابی 2، انتخاب و بازیابی:} آزمون نسبت درست‌نمایی؛
\lr{AIC}، \lr{AICc}، \lr{BIC} و وزن‌های \lr{Akaike}؛ بازیابی پارامتر و
مدل؛ \lr{identifiability} و \lr{trade-off} پارامترها & مطالعه‌ی بازیابی:
نمودار پارامتر واقعی در برابر بازیابی‌شده؛ \lr{confusion matrix} بازیابی
مدل \lr{(AIC/BIC)} & تفسیر معیارهای اطلاعاتی؛ طراحی مطالعه‌ی بازیابی؛
تشخیص \lr{non-identifiability} & \lr{F\lr{\&}L} فصل 10 & \textbf{تشکیل
گروه‌های پروژه (1 تا 2 نفر)؛ انتشار راهنمای پیشنهاد پروژه} \\
8 & 1\lr{--}3 & \textbf{مرور میان‌ترم:} یک مطالعه‌ی موردی یکپارچه از کل مسیر &
آزمون میان‌ترم (در زمان \lr{lab}، 90 دقیقه): سؤال مفهومی + خواندن و تفسیر ، فصل‌های 1 تا 3
& \lr{---} & مرور هفته‌های 1 تا 7 & \textbf{میان‌ترم} \\
9 & 4 & \textbf{مدل‌سازی بیزی 1:} پیشین، درست‌نمایی، پسین، شواهد
\lr{(evidence)}؛ \lr{conjugacy beta-binomial}؛ تقریب شبکه‌ای؛ خلاصه‌های
پسین و \lr{HDI}؛ \lr{credible interval} در برابر بازه‌ی اطمینان & پسین
شبکه‌ای برای احتمال یادآوری؛ شبکه‌ی دوبعدی برای \lr{SDT}؛ تغییر پیشین و
حجم نمونه؛ شبیه‌سازی پیش‌بینی پیشین & محاسبه‌ی پسین شبکه‌ای؛ توضیح نقش
پیشین؛ تفسیر \lr{credible interval} & \lr{K} فصل 5\lr{--}6؛ \lr{L\lr{\&}W} فصل 1 و
3؛ \lr{MKL} فصل 1؛ \lr{F\lr{\&}L} فصل 6؛ عظیمی‌فر و احمدی (برآورد بیزی) &
\textbf{موعد \lr{A3} · موعد پیشنهاد پروژه} \\
10 & 4 (+3) & \textbf{مدل‌سازی بیزی 2:} بررسی پیش‌بینی پسین \lr{(PPC)}؛
\lr{marginal likelihood} و \textbf{\lr{Bayes factor}}؛
\lr{Savage\lr{--}Dickey}؛ حساسیت به پیشین؛ \lr{BIC} به‌عنوان تقریب & \lr{PPC}
برای مدل فراموشی؛ \lr{Bayes factor Savage\lr{--}Dickey} برای \(d' > 0\)؛
\lr{Bayes factor} بین مدل‌های فراموشی با انتگرال‌گیری شبکه‌ای & انجام
\lr{PPC}؛ محاسبه و تفسیر \lr{Bayes factor}؛ ارزیابی حساسیت به پیشین &
\lr{L\lr{\&}W} فصل 7؛ \lr{F\lr{\&}L} فصل 11؛ \lr{K} فصل 10 & \textbf{انتشار
\lr{A4:}} تحلیل بیزی با شبکه (موعد: هفته‌ی 12) \\
11 & 5 & \textbf{\lr{Monte Carlo} و \lr{Metropolis\lr{--}Hastings:}}
انتگرال‌گیری \lr{Monte Carlo}؛ زنجیره‌ی مارکوف؛ \lr{MH}؛ تنظیم
\lr{proposal}، نرخ پذیرش، \lr{burn-in} & \textbf{\lr{MH} از صفر} برای
مدل فراموشی؛ \lr{trace plot}؛ مقایسه با پسین شبکه‌ای &
توضیح اینکه چرا \lr{MCMC }کار می‌کند؛ تحلیل \lr{MH}؛ تنظیم \lr{proposal}
& \lr{K} فصل 7؛ \lr{F\lr{\&}L} فصل 7 & \lr{---} \\
12 & 5 & \textbf{نمونه‌گیری \lr{Gibbs} و تشخیص همگرایی:} \lr{full
conditional}ها؛ \lr{Gibbs} برای مدل نرمال؛ \lr{MH-within-Gibbs}
(پیشرفته)؛ چند زنجیره، \textbf{\lr{R-hat}}، \lr{ESS}، خودهمبستگی &
\lr{Gibbs }از صفر ؛ تحلیل \lr{R-hat }و \lr{ESS }و مقایسه با ؛ خراب کردن عمدی همگرایی و تشخیص آن
&
تحلیل \lr{Gibbs}؛ تفسیر \lr{R-hat}، \lr{ESS }و \lr{trace plot}؛ قضاوت درباره‌ی قابل‌اعتماد بودن زنجیره
& \lr{K} فصل 7؛ \lr{MKL} فصل 2؛ \lr{F\lr{\&}L} فصل 7 & \textbf{موعد \lr{A4} ·
انتشار \lr{A5:}} \lr{MCMC} + تشخیص همگرایی (موعد: هفته‌ی 14) \\
13 & 5 &
احتمالاتی و مدل‌های سلسله‌مراتبی: ؛ \lr{HMC/NUTS }در سطح مفهومی (پیشرفته)؛ \lr{partial pooling }و \lr{shrinkage (}مرتبط با هفته‌ی 5
&
\lr{SDT }سلسله‌مراتبی در : تشخیص با ، \lr{PPC}، مقایسه با برازش‌های \lr{MLE }فردی؛ (پیشرفته) \lr{LOO/WAIC}
&
تعریف مدل در ؛ توضیح \lr{partial pooling}؛ گزارش یک گردش کار (\lr{workflow) }بیزی کامل
& \lr{L\lr{\&}W} فصل 11؛ \lr{K} فصل 9؛ \lr{F\lr{\&}L} فصل 9؛ \lr{MKL} فصل 2 و 10 &
\textbf{گزارش پیشرفت پروژه} (فقط بازخورد) \\
14 & 6 & \textbf{مدل‌سازی در علوم شناختی 1:} دسته‌بندی \lr{(prototype} در
برابر \lr{exemplar}؛ \textbf{\lr{GCM}})؛ انتخاب/زمان واکنش
(\textbf{\lr{drift-diffusion}}، \lr{EZ-diffusion)} & برازش \lr{GCM} با
\lr{MLE} روی داده‌ی ساختار دسته‌بندی 5-4؛ \lr{EZ-diffusion} روی زمان‌های
واکنش شبیه‌سازی‌شده؛ تفسیر پارامترها & برازش یک مدل \lr{exemplar}؛ نگاشت
پارامترهای \lr{DDM} به سازه‌های شناختی & \lr{L\lr{\&}W} فصل 17؛ \lr{F\lr{\&}L} فصل
14 {[}بررسی شود{]}؛ \lr{Verguts} (فصل‌های مرتبط {[}بررسی شود{]}) &
\textbf{موعد \lr{A5}} \\
15 & 6 & \textbf{مدل‌سازی در علوم شناختی 2:} \lr{Rescorla\lr{--}Wagner} و
\lr{Q-learning} با \lr{softmax}؛ \emph{مدل‌های بیزی شناخت} در برابر
\emph{تحلیل داده‌ی بیزی}؛ مدل‌های پیوندگرا و ارتباط با هوش مصنوعی & برازش
\lr{Q-learning} (\(\alpha\)، \(\beta\)) روی داده‌ی \lr{bandit}
شبیه‌سازی‌شده؛ بازیابی پارامتر؛ (پیشرفته) نسخه‌ی سلسله‌مراتبی در & برازش یک
مدل \lr{RL}؛ ارزیابی \lr{identifiability}؛ مقایسه‌ی شناخت بیزی با آمار
بیزی & \lr{Verguts} (فصل‌های یادگیری/\lr{RL} {[}بررسی شود{]})؛ \lr{F\lr{\&}L}
فصل 13 {[}بررسی شود{]}؛ \lr{Anderson} (فصل‌های مرتبط {[}بررسی شود{]}) &
کار روی پروژه‌ی پایانی \\
16 & \lr{---} & \textbf{ارائه‌ی پروژه‌های پایانی}؛ جمع‌بندی اصول مدل‌سازی خوب &
ادامه‌ی ارائه‌ها؛ بازخورد همتایان & ارائه و نقد یک مطالعه‌ی مدل‌سازی & \lr{---} &
ارائه‌ها · گزارش + پیش از آزمون پایانی \\
دوره‌ی امتحانات & 1\lr{--}6 &
آزمون پایان‌ترم (کتبی، جامع، با تأکید بر فصل‌های 4 تا 6؛ برگه‌ی فرمول داده می‌شود؛ خواندن و تفسیر ، بدون لپ‌تاپ
& \lr{---} & \lr{---} & \lr{---} & \textbf{آزمون پایان‌ترم} \\
\end{longtable}

\textbf{اختصارها:} \lr{F\lr{\&}L} = \lr{Farrell} \lr{\&} \lr{Lewandowsky} (2018) ·
\lr{L\lr{\&}W} = \lr{Lee} \lr{\&} \lr{Wagenmakers} (2014) · \lr{K} = \lr{Kruschke}
(2014) · \lr{MKL} = \lr{Martin, Kumar} \lr{\&} \lr{Lao} (2021). شماره‌ی فصل‌ها
از حافظه نوشته شده‌اند. آن‌هایی که {[}بررسی شود{]} دارند کمتر قطعی‌اند؛
لطفاً همه را با خود کتاب‌ها مقایسه کنید.

\subsection{نمره‌دهی (از 20)}

\begin{longtable}[]{@{}
  >{\raggedright\arraybackslash}p{(\linewidth - 2\tabcolsep) * \real{0.5000}}
  >{\raggedright\arraybackslash}p{(\linewidth - 2\tabcolsep) * \real{0.5000}}@{}}
\caption{\centering نمره‌دهی (از 20)}\tabularnewline
\toprule\noalign{}
\begin{minipage}[b]{\linewidth}\raggedright
بخش
\end{minipage} & \begin{minipage}[b]{\linewidth}\raggedright
نمره
\end{minipage} \\
\midrule\noalign{}
\endfirsthead
\toprule\noalign{}
\begin{minipage}[b]{\linewidth}\raggedright
بخش
\end{minipage} & \begin{minipage}[b]{\linewidth}\raggedright
نمره
\end{minipage} \\
\midrule\noalign{}
\endhead
\bottomrule\noalign{}
\endlastfoot
تأیید تمرین‌های \lr{lab} (تمرین‌های هفتگی سر کلاس) & 2.0 \\
تکالیف \lr{A1 }تا \lr{A5 (}هرام 1.0 & 5.0 \\
کوییزهای کوتاه سر کلاس (حدود 10 کوییز، 2 نمره‌ی پایین حذف) & 1.0 \\
آزمون میان‌ترم (هفته‌ی 8) & 3.0 \\
پروژه‌ی پایانی: پیشنهاد 1.0 · ارائه 2.0 · گزارش + بازتولیدپذیر 3.0 &
6.0 \\
آزمون پایان‌ترم (جامع) & 3.0 \\
\textbf{جمع} & \textbf{20.0} \\
\end{longtable}

\subsection{پروژه‌ی پایانی: نمای کلی}

\begin{itemize}
\tightlist
\item
  \textbf{گروه‌های 1 یا 2 نفره.} حجم کار با تعداد اعضا بیشتر می‌شود. همه
  حداقل دو مدل رقیب را برازش می‌کنند، یک مقایسه‌ی مدل و یک تحلیل بازیابی
  پارامتر انجام می‌دهند. \textbf{گروه‌های دونفره علاوه بر این} یک تحلیل
  بیزی/\lr{MCMC} (یا یک مدل سلسله‌مراتبی) انجام می‌دهند و بیانیه‌ی سهم فردی
  تحویل می‌دهند.
\item
  \textbf{داده:} داده‌ی شبیه‌سازی‌شده، یک مجموعه‌داده‌ی باز، یا داده‌ی آزمایشی
  خود دانشجو.
\end{itemize}
